% =====================================================================
\section{Triac y control de fase}
\label{sec:triac}
% =====================================================================

\subsection{El triac: un tiristor bidireccional}

El triac (\emph{triode for alternating current}) integra en una sola pastilla el equivalente de
\textbf{dos tiristores en antiparalelo con una compuerta común} \diapo{6}. Sus terminales principales
no se llaman ánodo y cátodo sino MT1 y MT2 (\emph{main terminals}); la compuerta G está del lado de MT1
y todas las tensiones se miden respecto de MT1. Es una estructura de cinco regiones ($npnpn$ vista en
un corte): según el signo de $V_{MT2-MT1}$ conduce una u otra mitad (Pierret §13.3).

\begin{figure}[H]
\centering
\begin{minipage}{0.3\textwidth}\centering
\begin{circuitikz}
  \draw (0,0) node[below] {MT1} to[triac,n=T] (0,2.6) node[above] {MT2};
  \draw (T.gate) -- ++(-0.6,0) node[left] {G};
\end{circuitikz}
\end{minipage}
\begin{minipage}{0.3\textwidth}\centering
\begin{circuitikz}
  \draw (0,0) node[below] {MT1} -- (0,0.3) -- (-0.8,0.3) to[thyristor,n=S1] (-0.8,2.3) -- (0,2.3) -- (0,2.6) node[above] {MT2};
  \draw (0,0.3) -- (0.8,0.3) to[thyristor,invert,n=S2] (0.8,2.3) -- (0,2.3);
\end{circuitikz}
\end{minipage}
\begin{minipage}{0.36\textwidth}\centering
\begin{tikzpicture}[x=1cm,y=1cm]
  \fill[azul!30] (0,2.2) rectangle (1.1,2.6);
  \fill[rojo!15] (1.1,2.2) rectangle (2.6,2.6);
  \fill[rojo!15] (0,1.8) rectangle (2.6,2.2);
  \fill[azul!8] (0,0.8) rectangle (2.6,1.8);
  \fill[rojo!15] (0,0.4) rectangle (2.6,0.8);
  \fill[rojo!15] (0,0) rectangle (1.5,0.4);
  \fill[azul!30] (1.5,0) rectangle (2.6,0.4);
  \fill[azul!30] (0,0.4) rectangle (0.5,0.62);
  \draw (0,0) rectangle (2.6,2.6);
  \node[font=\scriptsize] at (0.55,2.4) {$n$}; \node[font=\scriptsize] at (1.85,2.4) {$p$};
  \node[font=\scriptsize] at (1.3,2.0) {$p$}; \node[font=\scriptsize] at (1.3,1.3) {$n^-$};
  \node[font=\scriptsize] at (1.3,0.6) {$p$}; \node[font=\scriptsize] at (0.75,0.2) {$p$}; \node[font=\scriptsize] at (2.05,0.2) {$n$};
  \node[font=\scriptsize] at (0.25,0.51) {$n$};
  \fill[gris] (0,2.6) rectangle (2.6,2.7); \fill[gris] (0.6,-0.1) rectangle (2.6,0);
  \fill[gris] (0,-0.1) rectangle (0.4,0);
  \node[font=\small] at (1.3,2.95) {MT2}; \node[font=\small] at (1.6,-0.35) {MT1}; \node[font=\small] at (0.2,-0.35) {G};
\end{tikzpicture}
\end{minipage}
\caption{Triac: símbolo, equivalente con dos tiristores en antiparalelo y corte esquemático
\diapo{6}. Las regiones $n$ junto a cada terminal hacen de cátodo de uno u otro tiristor según la
polaridad; la $n$ junto a la compuerta permite dispararlo con $\IG$ de ambos signos.}
\label{fig:triac}
\end{figure}

\begin{figure}[H]
\centering
\begin{tikzpicture}
\begin{axis}[
  width=0.78\textwidth, height=6.4cm,
  axis lines=middle, xmin=-7, xmax=7.5, ymin=-4.2, ymax=4.2,
  xtick=\empty, ytick=\empty,
  xlabel={$V_{MT2-MT1}$}, ylabel={$I$},
  every axis plot/.append style={line width=1.1pt},
]
  \addplot[azul,smooth] coordinates {(0.65,1.1) (0.85,1.7) (1.1,2.8) (1.4,4.0)};
  \addplot[rojo,smooth] coordinates {(0,0) (2,0.15) (4,0.3) (5.5,0.55) (5.9,1.0)};
  \addplot[rojo,dashed,smooth] coordinates {(5.9,1.0) (3,1.2) (0.75,1.3)};
  \addplot[naranja,dashed,smooth] coordinates {(3.0,0.6) (2,1.1) (0.75,1.3)};
  \addplot[naranja,smooth] coordinates {(0,0) (2.6,0.3) (3.0,0.6)};
  \addplot[azul,smooth] coordinates {(-0.65,-1.1) (-0.85,-1.7) (-1.1,-2.8) (-1.4,-4.0)};
  \addplot[rojo,smooth] coordinates {(0,0) (-2,-0.15) (-4,-0.3) (-5.5,-0.55) (-5.9,-1.0)};
  \addplot[rojo,dashed,smooth] coordinates {(-5.9,-1.0) (-3,-1.2) (-0.75,-1.3)};
  \addplot[naranja,dashed,smooth] coordinates {(-3.0,-0.6) (-2,-1.1) (-0.75,-1.3)};
  \addplot[naranja,smooth] coordinates {(0,0) (-2.6,-0.3) (-3.0,-0.6)};
  \node[font=\small,anchor=south west] at (axis cs:5.9,1.05) {$V_{BO}$};
  \node[font=\small,anchor=east] at (axis cs:-5.95,-1.0) {$-V_{BO}$};
  \node[font=\small] at (axis cs:3.2,3.3) {cuadrante I};
  \node[font=\small] at (axis cs:-3.2,-3.3) {cuadrante III};
\end{axis}
\end{tikzpicture}
\caption{Característica del triac \diapo{6}: la del tiristor repetida en el tercer cuadrante. Se
dispara en los dos sentidos, con corriente de compuerta entrante o saliente.}
\label{fig:triac-iv}
\end{figure}

\begin{table}[H]
\centering\small
\begin{tabular}{cccl}
\toprule
Modo & $V_{MT2-MT1}$ & $\IG$ & Sensibilidad \\
\midrule
I$^+$   & $+$ & $+$ (entra) & la mejor \\
I$^-$   & $+$ & $-$ (sale)  & buena \\
III$^-$ & $-$ & $-$ (sale)  & buena \\
III$^+$ & $-$ & $+$ (entra) & mala (se evita) \\
\bottomrule
\end{tabular}
\caption{Modos de disparo del triac. En la práctica se usan I$^-$ y III$^-$ (compuerta siempre
negativa) o I$^+$ y III$^-$ (pulso del mismo signo que $V_{MT2}$).}
\label{tab:triac-cuadrantes}
\end{table}

\begin{intuicion}[Tiristor o triac \diapo{7}]
El triac es cómodo (un solo componente y un solo disparo para alterna), pero es más lento y soporta
menos $\dd v/\dd t$ de conmutación: al pasar la corriente por cero, la carga que quedó de un
semiciclo puede reencender el dispositivo en el semiciclo siguiente. Por eso se usa en \textbf{baja
potencia} (atenuadores de luz, control de motores universales, relés de estado sólido) y el
\textbf{tiristor en media, alta y muy alta potencia} (rectificadores controlados, transmisión HVDC,
arrancadores).
\end{intuicion}

\subsection{Control de fase}

La aplicación clásica de ambos es el \textbf{control de fase}: en cada semiciclo de la red se
dispara el dispositivo con un retardo $\alpha$ (ángulo de disparo) respecto del cruce por cero. El
dispositivo conduce desde $\omega t=\alpha$ hasta que la corriente pasa por cero (con carga
resistiva, en $\omega t=\pi$) y la carga recibe una ``porción'' de cada semiciclo \diapo{7}.

\begin{figure}[H]
\centering
\begin{circuitikz}
  \draw (0,0) to[sV,l={$v_s=V_m\sin\omega t$}] (0,2.6) -- (1.2,2.6) to[thyristor,n=S] (3.6,2.6)
        -- (5.0,2.6) to[R,l=$R$] (5.0,0) -- (0,0);
  \node[draw,font=\small,align=center] (ctl) at (2.9,3.9) {control de compuerta};
  \draw (S.gate) -- (S.gate |- ctl.south);
  \node at (5.9,2.3) {$+$}; \node at (5.9,1.3) {$v_o$}; \node at (5.9,0.3) {$-$};
\end{circuitikz}

\medskip
\begin{tikzpicture}
\begin{axis}[
  width=0.9\textwidth, height=4.0cm, name=a1,
  xmin=0, xmax=720, ymin=-1.15, ymax=1.15,
  xtick={0,60,180,360,420,540,720}, xticklabels={0,$\alpha$,$\pi$,$2\pi$,,$3\pi$,$4\pi$},
  ytick=\empty, ylabel={tiristor},
  tick label style={font=\scriptsize}, samples=361, domain=0:720,
]
  \addplot[gris,dashed] {sin(x)};
  \addplot[rojo,line width=1.1pt] {(mod(x,360)>=60 && mod(x,360)<=180) ? sin(x) : 0};
  \addplot[violeta,ycomb,line width=1.2pt] coordinates {(60,0.35) (420,0.35)};
\end{axis}
\begin{axis}[
  width=0.9\textwidth, height=4.0cm, at={(a1.south)}, anchor=north, yshift=-0.2cm,
  xmin=0, xmax=720, ymin=-1.15, ymax=1.15,
  xtick={0,60,180,240,360,420,540,600,720}, xticklabels={0,$\alpha$,$\pi$,$\pi{+}\alpha$,$2\pi$,,$3\pi$,,$4\pi$},
  ytick=\empty, ylabel={triac},
  tick label style={font=\scriptsize}, samples=361, domain=0:720,
  xlabel={$\omega t$},
]
  \addplot[gris,dashed] {sin(x)};
  \addplot[verde!60!black,line width=1.1pt] {(mod(x,180)>=60) ? sin(x) : 0};
  \addplot[violeta,ycomb,line width=1.2pt] coordinates {(60,0.35) (240,-0.35) (420,0.35) (600,-0.35)};
\end{axis}
\end{tikzpicture}
\caption{Control de fase con carga resistiva, $\alpha=\ang{60}$ como en la diapo~7. Arriba: tiristor
(media onda); los pulsos (violeta) en el semiciclo negativo no tienen efecto porque el tiristor está
en inversa. Abajo: triac (onda completa), disparado en los dos semiciclos.}
\label{fig:control-fase}
\end{figure}

Integrando la porción de senoide entre $\alpha$ y $\pi$:

\begin{clave}[Control de fase con carga resistiva]
\[
  \text{Tiristor (media onda):}\qquad
  V_{\mathrm{med}}=\frac{1}{2\pi}\int_\alpha^\pi V_m\sin\theta\dd\theta=\frac{V_m}{2\pi}(1+\cos\alpha)
\]
\[
  \text{Triac (onda completa):}\qquad
  V_{\mathrm{ef}}=\sqrt{\frac1\pi\int_\alpha^\pi V_m^2\sin^2\theta\dd\theta}
  =\frac{V_m}{\sqrt2}\sqrt{1-\frac{\alpha}{\pi}+\frac{\sin2\alpha}{2\pi}},
  \qquad P=\frac{V_{\mathrm{ef}}^2}{R}.
\]
\end{clave}

\begin{ejemplo}[Atenuador de \SI{220}{V} con triac]
Lámpara de $R=\SI{100}{\ohm}$ (supuesta constante) en la red de \SI{220}{V} ($V_m=\SI{311}{V}$):
\begin{center}\small
\begin{tabular}{lcccc}
\toprule
$\alpha$ & \ang{0} & \ang{60} & \ang{90} & \ang{120} \\
\midrule
$V_{\mathrm{ef}}$ (triac) & \SI{220}{V} & \SI{197}{V} & \SI{156}{V} & \SI{97}{V} \\
$P$ & \SI{484}{W} & \SI{389}{W} & \SI{242}{W} & \SI{95}{W} \\
$V_{\mathrm{med}}$ (tiristor, media onda) & \SI{99}{V} & \SI{74}{V} & \SI{50}{V} & \SI{25}{V} \\
\bottomrule
\end{tabular}
\end{center}
Con $\alpha=\ang{90}$ la potencia es exactamente la mitad: se entrega medio semiciclo de cada
uno, y por simetría de $\sin^2$ cada mitad tiene la misma energía.
\end{ejemplo}

\begin{intuicion}[El precio del control de fase]
La corriente de línea deja de ser senoidal: el encendido brusco en $\alpha$ genera armónicos e
interferencia (de ahí los filtros LC de los atenuadores) y el factor de potencia empeora aun con
carga resistiva. Para generar el pulso en $\alpha$, los atenuadores simples usan una red RC que
carga un capacitor hasta la tensión de disparo de un \emph{diac} (dos junturas en antiparalelo que
rompen a $\sim\SI{30}{V}$).
\end{intuicion}
